\documentclass[11pt,a4paper]{article}

% --- Core LaTeX Packages ---
\usepackage[utf8]{inputenc}
\usepackage[margin=1in]{geometry}
\usepackage{amsmath,amssymb,amsthm}
\usepackage{xcolor}
\usepackage{listings}
\usepackage{hyperref}
\usepackage{tabularx}
\usepackage{graphicx}

% --- Unified Color Palette ---
\definecolor{primaryblue}{RGB}{24, 75, 140}
\definecolor{codebg}{RGB}{248, 249, 250}
\definecolor{codeframe}{RGB}{210, 215, 222}
\definecolor{codegreen}{RGB}{40, 130, 60}
\definecolor{codeblue}{RGB}{20, 90, 180}
\definecolor{codegray}{RGB}{120, 120, 120}

% --- Hyperref Setup ---
\hypersetup{
    colorlinks=true,
    linkcolor=primaryblue,
    citecolor=primaryblue,
    urlcolor=primaryblue,
    pdftitle={Dynamic Programming},
    pdfauthor={Antigravity Engineering}
}

% --- Theorem & Definition Environments ---
\theoremstyle{definition}
\newtheorem{definition}{Definition}[section]
\newtheorem{theorem}{Theorem}[section]
\newtheorem{lemma}[theorem]{Lemma}
\newtheorem{example}{Example}[section]

% --- Callout Box Environment ---
\newenvironment{calloutbox}[1]{
    \par\vspace{0.3cm}
    \noindent\begin{tabular}{|p{0.96\textwidth}|}
    \hline
    \vspace{0.1cm}
    \textbf{#1}\par\vspace{0.1cm}
}{
    \vspace{0.1cm}
    \\ \hline
    \end{tabular}
    \vspace{0.3cm}\par
}

% --- Modern C++ Code Listing Setup ---
\lstset{
    language=C++,
    basicstyle=\footnotesize\ttfamily,
    keywordstyle=\color{codeblue}\bfseries,
    commentstyle=\color{codegreen}\itshape,
    stringstyle=\color{orange!80!black},
    numberstyle=\tiny\color{codegray},
    numbers=left,
    numbersep=8pt,
    backgroundcolor=\color{codebg},
    frame=single,
    rulecolor=\color{codeframe},
    breaklines=true,
    breakatwhitespace=true,
    tabsize=4,
    showstringspaces=false,
    captionpos=b
}

% --- Title Block ---
\title{\vspace{-1.5cm}\textbf{\Huge Dynamic Programming}\\[0.3cm]
\Large States, Transitions, Memoization, Tabulation, and Classic Table-Filling Patterns}
\author{Technical Monograph}
\date{\today}

\begin{document}

\maketitle

\begin{abstract}
\noindent
Dynamic programming is a method for solving complex problems by breaking them down into simpler, recurring overlapping subproblems. This document guides you from basic recursive approaches to efficient tabular techniques using caching and optimal substructure. Designed for programmers familiar with C++ and basic recursion, it covers everything from 1D sequences to advanced multidimensional and knapsack patterns.
\end{abstract}

\tableofcontents
\newpage

\section{Foundations}

If you know recursion, you already know how to solve problems by breaking them down into smaller pieces. However, pure recursion often recalculates the same results multiple times. Dynamic Programming (DP) is the art of identifying this repeated work and saving the results so you only compute them once.

\subsection{What is a Subproblem?}

A subproblem is simply a smaller version of your original problem. For example, if you want to find the shortest path from your house to the grocery store, a subproblem might be finding the shortest path from the intersection halfway there to the store.

\subsection{Repeated Work in Plain Recursion}

Consider the classic Fibonacci sequence, where each number is the sum of the two preceding ones. The base cases are $f(0) = 0$ and $f(1) = 1$. Let us trace the recursive calls for $f(5)$ in a naive recursive approach:

Notice how $f(3)$ is computed entirely from scratch twice. $f(2)$ is computed three times. The number of calls grows exponentially. The total number of calls to compute $f(n)$ using this naive branching is exactly $2 \times f(n+1) - 1$. For $f(5)$, where $f(6) = 8$, the formula yields $2 \times 8 - 1 = 15$ calls, exactly matching the 15 nodes in the tree above. By the time you ask for $f(40)$, plain recursion makes over 330 million calls!

By remembering (caching) what we compute, we trim the tree. The first time a node is evaluated, we save the result. Subsequent requests for the same value are immediate cache hits:

With memoization, the call count for $f(5)$ drops to just 9 nodes ($2n - 1$). For $f(40)$, it takes a mere 79 calls.

\subsection{The Two Properties of DP}

A problem is solvable with DP if it has two properties:
\begin{enumerate}
    \item \textbf{Overlapping Subproblems}: You encounter the exact same subproblem multiple times, as seen with $f(3)$ appearing multiple times above. In programming, if a recursive function takes the same parameters repeatedly, it has overlapping subproblems.
    \item \textbf{Optimal Substructure}: The optimal solution to a large problem can be constructed from the optimal solutions of its subproblems. Everyday analogy: the shortest path from City A to City C going through City B is the shortest path from A to B combined with the shortest path from B to C.
\end{enumerate}

When you have these two properties, you can apply DP. At its core, DP is just ``recursion plus a memory cache'' (memoization) or ``filling a table in the right order'' (tabulation). You map the inputs of your recursive function to a spot in an array or a \texttt{std::unordered\_map}, and check if you already solved it before doing the work.

\subsection{Core Vocabulary}

\begin{itemize}
    \item \textbf{State}: A set of variables that perfectly describes a specific subproblem. In Fibonacci, the state is just the integer $n$.
    \item \textbf{Transition} (or Recurrence): The formula that expresses the answer for a state in terms of smaller states. For Fibonacci: $f(n) = f(n-1) + f(n-2)$.
    \item \textbf{Base Case}: The smallest states that do not need recursion to solve. For Fibonacci: $f(0) = 0$ and $f(1) = 1$.
    \item \textbf{Answer State}: The specific state that holds the final answer to the overall problem. For Fibonacci to the $n$-th number, it is $f(n)$.
    \item \textbf{Evaluation Order}: The sequence in which you must solve subproblems so that whenever you need to compute a transition, the smaller states are already solved. For Fibonacci, you evaluate from smallest $n$ to largest $n$.
\end{itemize}

\begin{calloutbox}{Key Insight: DP is just a Directed Acyclic Graph}
Every DP problem is a graph. The nodes are the states. The directed edges are the dependencies (transitions). The evaluation order is a topological sort of this graph, guaranteeing you never need an answer that has not been computed yet.
\end{calloutbox}

\section{Main Topics}

\subsection{Top-down Memoization vs Bottom-up Tabulation}

There are two primary ways to write a DP solution. Both have the same time complexity, but they differ in how they compute the states.

Let us use a simple example: counting the number of ways to climb $n$ stairs if you can take 1 or 2 steps at a time. The transition is $W(n) = W(n-1) + W(n-2)$. If we represent this as a state dependency graph, we see that Step 3 needs Step 1 and Step 2:

\textbf{Top-down Memoization} starts at the final goal $W(n)$ and recursively asks for smaller answers. Whenever it computes an answer, it saves it in a cache (like an array initialized to -1). Before computing a state, it checks if the cache already has the answer.

\textbf{Bottom-up Tabulation} starts at the base cases and iteratively fills an array up to the final answer. It uses a loop instead of recursion.

\begin{calloutbox}{Decision Guide: Memoization vs Tabulation}
\begin{itemize}
    \item Use \textbf{Tabulation} (Bottom-up) by default. It avoids function call overhead, never hits recursion depth limits, and makes space optimization easier.
    \item Use \textbf{Memoization} (Top-down) when the state space is very sparse (meaning you only need to visit a small fraction of all possible states), or if the transitions are extremely complex to figure out via loops.
\end{itemize}
\end{calloutbox}

\subsection{The 5-Step Recipe}

To reliably solve DP problems, always follow these 5 steps in order:
\begin{enumerate}
    \item \textbf{Define the state in words}: What exactly does \texttt{dp[i]} represent?
    \item \textbf{Write the transition}: How do you compute \texttt{dp[i]} from previous states?
    \item \textbf{Set base cases}: What are the simplest states?
    \item \textbf{Choose evaluation order}: How should your loops run to respect dependencies?
    \item \textbf{Locate the answer}: Which \texttt{dp} cell holds the final answer?
\end{enumerate}

We will apply this to the \textbf{House Robber} problem in the Worked Examples section. The core idea is that you want to maximize money without choosing adjacent houses, leading to a transition where you either skip the current house or rob it.

\subsection{1D DP: Coins and Subsequences}

1D DP problems have a state defined by a single index. A classic example is the \textbf{Longest Increasing Subsequence (LIS)} problem, where the state \texttt{dp[i]} tracks the best sequence ending exactly at index $i$.

Another fundamental 1D problem is \textbf{Coin Change - Counting Ways}. The loop order matters deeply when counting combinations versus permutations. Let us use coins $\{1, 2\}$ and target amount $3$:

If the outer loop is the coins and the inner loop is the amounts, you count \textit{combinations} (order does not matter, for example, $1+2$ is the same as $2+1$). The table evolves coin by coin:

If the outer loop is the amount and the inner loop is the coins, you count \textit{permutations} (order matters). The table evolves amount by amount:

\subsection{2D DP on Grids}

Many problems give you a 2D grid and ask you to navigate from the top-left to the bottom-right, only moving right or down. The state \texttt{dp[i][j]} represents the minimum cost or total paths to reach row $i$, col $j$. The transition checks the cell above and the cell to the left.

\subsection{Contiguous Subarrays and Substrings (The `Ends Exactly At'' Pattern)}

Many dynamic programming problems deal with substrings or contiguous subarrays (unlike subsequences, where elements can be skipped). Because a substring must remain perfectly unbroken, tracking the `best answer seen so far'' in \texttt{dp[i]} breaks the DP transition---if you only know the best overall score, you have no way to attach the current character $i$ to it, because the best sequence might have ended several indices ago.

The universal solution to this is the `Ends Exactly At'' pattern. We rigidly define our state:
\begin{calloutbox}{The Substring DP State Rule}
\texttt{dp[i]} must represent the optimal answer for a valid sequence that \textbf{stops exactly at index $i$}.
\end{calloutbox}

By forcing the sequence to end exactly at the current character, we guarantee that any sequence extending it will remain unbroken. 

\textbf{Worked Trace: Kadane's Algorithm for Maximum Subarray Sum} \\
Given the array [-2, 1, -3, 4, -1, 2]. We define \texttt{dp[i]} as the maximum sum of a contiguous subarray ending \textit{exactly} at $i$. At each step, we have two choices: extend the subarray ending at $i-1$, or start a brand new subarray at $i$.
\begin{enumerate}
    \item $i=0$ (val -2): \texttt{dp[0]} = -2. 
    \item $i=1$ (val 1): Extend (-2 + 1 = -1) or Start fresh (1). Max is 1. \texttt{dp[1]} = 1.
    \item $i=2$ (val -3): Extend (1 + -3 = -2) or Start fresh (-3). Max is -2. \texttt{dp[2]} = -2.
    \item $i=3$ (val 4): Extend (-2 + 4 = 2) or Start fresh (4). Max is 4. \texttt{dp[3]} = 4.
    \item $i=4$ (val -1): Extend (4 + -1 = 3) or Start fresh (-1). Max is 3. \texttt{dp[4]} = 3.
    \item $i=5$ (val 2): Extend (3 + 2 = 5) or Start fresh (2). Max is 5. \texttt{dp[5]} = 5.
\end{enumerate}
Because \texttt{dp[i]} only tracks the sequence ending at $i$, the global optimal answer is simply the maximum value found anywhere in the \texttt{dp} array (which is 5). This identical `ends exactly at'' logic applies to Longest Valid Parentheses, Longest Palindromic Substring, and any contiguous-segment problem.

\subsection{Two-Sequence DP: Longest Common Subsequence and Edit Distance}

When comparing two strings or sequences, the state needs two indices: $i$ for the first string and $j$ for the second.
In \textbf{Edit Distance}, we find the minimum insert, delete, or replace operations to transform one string into another. In \textbf{LCS}, we find the longest shared characters in order. Both rely on checking if the current characters match, and branching accordingly.

\subsection{The Knapsack Family}

\textbf{0/1 Knapsack}: You have items with weights and values. You have a bag of capacity $W$. Maximize value without exceeding capacity. You can take each item at most once. The state \texttt{dp[i][w]} tracks the best value using the first $i$ items and $w$ capacity.
When optimizing the space to a 1D array \texttt{dp[w]}, you must iterate the capacity \textit{backwards}. If you iterate forwards, you risk using the same item multiple times, turning it into the Unbounded Knapsack problem.

\textbf{Subset-Sum}: A boolean variant of knapsack. Can we find a subset that sums exactly to a target? The state is boolean.

\subsection{Interval DP}

In Interval DP, you combine adjacent elements (like multiplying matrices or merging piles). The state is defined by a left and right boundary: \texttt{dp[L][R]}. Because computing a large interval requires the answers to smaller sub-intervals, the evaluation order mandates looping over the \textit{length} of the interval first, starting from 1 up to $N$.

\subsection{Space Optimization (Rolling Arrays)}

Often, a state only depends on the immediate previous row. In a 2D grid, \texttt{dp[i][j]} only reads from row $i$ and row $i-1$. Once row $i$ is finished, row $i-2$ is completely useless. We can reduce an $O(N \times M)$ space complexity to $O(M)$ by keeping only two rows, swapping them as we progress.

\subsection{Bitmask DP}

When you have a small set of items ($N \le 20$), you can represent which items are visited using an integer where each bit is a boolean flag. The state is often \texttt{dp[mask][last\_visited]}. The transitions use bitwise operations (\texttt{<<}, \texttt{\&}, \texttt{|}) to check or set visited items.

\subsection{Reconstructing the Solution}

Sometimes you need the actual path, not just the max value. You can deduce the choice backwards by checking which transition gave the optimal value, or you can maintain an explicit \texttt{parent} array during the computation that stores the exact choice made at each state, allowing you to walk the pointers back to the start.

\subsection{DP vs Greedy vs Backtracking}

\begin{calloutbox}{Decision Guide: Backtracking vs Greedy vs DP}
\begin{itemize}
    \item \textbf{Backtracking}: Explores all possible paths. Used when you need to list EVERY valid configuration, or $N$ is tiny (for example, N=10).
    \item \textbf{Greedy}: Makes the locally optimal choice at every step. Extremely fast but often wrong! Example where Greedy fails: Coins \{1, 3, 4\}, amount 6. Greedy picks 4, then 1, 1 (3 coins). DP correctly picks 3, 3 (2 coins).
    \item \textbf{DP}: Explores all paths intelligently by merging overlapping subproblems. Used when you need an optimal value (min/max) or count, and subproblems overlap.
\end{itemize}
\end{calloutbox}

\begin{calloutbox}{When DP is the wrong tool}
If the problem asks you to find the lexicographically smallest path, or a path without returning to visited nodes (and $N$ is large), the state requires remembering the entire history. This destroys overlapping subproblems. DP only works when the past can be summarized by a small state.
\end{calloutbox}

\subsection{Counting Problems with a Modulus}

Because counting paths can yield huge numbers, problems often ask for the answer modulo a prime (like $10^9 + 7$). Apply the modulo operation at every single addition: \texttt{dp[i] = (dp[i] + dp[i-1]) \% MOD}. 

If you ever subtract, you must avoid negative numbers in C++: \texttt{dp[i] = (dp[i] - val + MOD) \% MOD}. For example, if doing calculations modulo 7, and you subtract 5 from 3: $(3 - 5) \dots -2 \dots (-2 + 7) \dots 5$. The \texttt{+ MOD} guarantees a positive result before the final modulo.

\begin{calloutbox}{Integer Overflow Without Modulo}
If a counting problem does not specify a modulus, the answer might still exceed a 32-bit integer. Always use \texttt{long long} (64-bit integers) for combinations or path counts unless bounds explicitly fit in a standard integer.
\end{calloutbox}

\begin{calloutbox}{Common Pitfalls}
\begin{itemize}
    \item \textbf{Off-by-one in base cases}: Forgetting to set $dp[0] = 1$ in counting problems, resulting in an array of zeroes.
    \item \textbf{Wrong sentinel values}: Initializing a minimum-cost array with exactly the maximum integer, which overflows and wraps to a negative number when you add a transition cost.
    \item \textbf{Wrong loop direction}: Overwriting a 1D array value before reading the older value needed for the current transition.
\end{itemize}
\end{calloutbox}

\begin{calloutbox}{How to find the state when stuck}
Ask yourself: \textit{What is the minimum amount of information I need to know about the past to make the correct decision right now?} Strip away anything you do not need. If you do not need the exact sequence of items picked, just their sum, the state is the sum.
\end{calloutbox}

\section{Key Algorithmic Patterns}

Why is naive recursion exponential? Because it explores a tree where many branches are identical. As shown earlier, $f(40)$ requires hundreds of millions of calls, but there are only 41 distinct states! The efficient technique is to identify the state space, initialize a data structure (the table or cache) to store the result of each state, and compute them exactly once.

\textbf{Where does this come from?}
This is valid because of the optimal substructure property: a state fully determines the future. If you know you are at index 5 with 10 capacity remaining, it does not matter \textit{how} you got there. The best decisions from this point onward are always the same. Storing the answer is safe.

Here are the canonical DP patterns you will see again and again:

\begin{itemize}
    \item \textbf{Linear Sequence (1D)}: State is just the index \texttt{i}. Usually represents a prefix array. Example trace: Array $[3, 1, 4]$. \texttt{dp[0]=3}, \texttt{dp[1]=4}, \texttt{dp[2]=8}.
    \item \textbf{Two-Index (2D Grid)}: State is \texttt{dp[i][j]}. You move right/down on a grid.
    \item \textbf{Two-Sequence}: State is \texttt{dp[i][j]}, consuming characters from two strings.
    \item \textbf{Knapsack}: State is \texttt{dp[item][capacity]}. You decide whether to take an item, branching to \texttt{w} or \texttt{w - weight}.
    \item \textbf{Interval}: State is \texttt{dp[L][R]}. You solve for small ranges (length 1, then length 2) and expand outwards to the full string.
    \item \textbf{Subset / Bitmask}: State is an integer representing a set of booleans. You iterate over bits to transition.
\end{itemize}

\section{Worked Examples}

\textbf{Example 1: The House Robber (1D Linear)}
Problem: Maximize money without stealing from adjacent houses. Values: $[2, 7, 9, 3, 1]$.
State: \texttt{dp[i]} = max money using prefix up to $i$.
Transition: $\texttt{dp[i]} = \max(\texttt{dp[i-1]} \text{ (skip)}, \texttt{dp[i-2]} + \text{val}[i] \text{ (rob)})$.
Trace:
\begin{enumerate}
    \item Base cases: \texttt{dp[0]} = 2. \texttt{dp[1]} = $\max(2, 7) = 7$.
    \item $i=2$ (val 9): skip=7, rob=$2+9=11$. \texttt{dp[2]} = 11.
    \item $i=3$ (val 3): skip=11, rob=$7+3=10$. \texttt{dp[3]} = 11.
    \item $i=4$ (val 1): skip=11, rob=$11+1=12$. \texttt{dp[4]} = 12.
\end{enumerate}

The final answer is 12.

\textbf{Example 2: Minimum Path Sum (2D Grid)}
Problem: Minimum cost to reach bottom-right.
Values:

State: \texttt{dp[i][j]} = $\text{grid}[i][j] + \min(\texttt{dp[i-1][j]}, \texttt{dp[i][j-1]})$.
Trace with arrows showing dependencies (cells read from left and above):

For cell (1,1) containing 5: It reads left (2) and above (4). $\min(2, 4) + 5 = 7$.
For cell (1,2) containing 1: It reads left (7) and above (5). $\min(7, 5) + 1 = 6$. The final answer is 6.

\textbf{Example 3: Longest Increasing Subsequence}

Problem: Find LIS on $[10, 9, 2, 5, 3, 7, 101, 18]$.
State: \texttt{dp[i]} is the LIS ending at index $i$.
Initialize all to 1.
\begin{enumerate}
    \item $i=0$ (10): \texttt{dp[0]} = 1.
    \item $i=1$ (9): No previous element is smaller. \texttt{dp[1]} = 1.
    \item $i=2$ (2): \texttt{dp[2]} = 1.
    \item $i=3$ (5): 2 is smaller. \texttt{dp[3]} = $\max(1, \texttt{dp[2]}+1) = 2$.
    \item $i=4$ (3): 2 is smaller. \texttt{dp[4]} = 2.
    \item $i=5$ (7): 2, 5, 3 are smaller. Max is from 5 or 3. \texttt{dp[5]} = 3.
    \item $i=6$ (101): All are smaller. Max is from 7. \texttt{dp[6]} = 4.
    \item $i=7$ (18): 10, 9, 2, 5, 3, 7 are smaller. Max is from 7. \texttt{dp[7]} = 4.
\end{enumerate}

The max value in \texttt{dp} is 4.

\textbf{Example 4: Edit Distance (Two-Sequence)}
Problem: Transform ``kitten'' to ``sitting''.
State: \texttt{dp[i][j]} = min edits for prefix $i$ of word1 and prefix $j$ of word2.
Trace of the full $7 \times 8$ table. If characters match, cost is \texttt{dp[i-1][j-1]}. Otherwise, $1 + \min(\text{insert}, \text{delete}, \text{replace})$.

Walk-back path to reconstruct:
\begin{enumerate}
    \item Start at $(6,7)$ with value 3. Characters \texttt{'n'} vs \texttt{'g'} differ. Minimum neighbor is $(5,6)$ with value 2 (Replace \texttt{'n'} with \texttt{'g'}).
    \item At $(5,6)$ with value 2. Characters \texttt{'e'} vs \texttt{'n'} differ. Minimum neighbor is $(4,5)$ with value 2 (Insert \texttt{'n'}).
    \item At $(4,5)$ with value 2. Characters \texttt{'t'} vs \texttt{'i'} differ. Minimum neighbor is $(3,4)$ with value 1 (Replace \texttt{'t'} with \texttt{'i'}).
    \item At $(3,4)$ with value 1. Characters \texttt{'t'} vs \texttt{'t'} match. Move diagonally to $(2,3)$ with value 1.
    \item At $(2,3)$ with value 1. Characters \texttt{'i'} vs \texttt{'t'} differ. Minimum neighbor is $(2,2)$ with value 1 (Insert \texttt{'t'}).
\end{enumerate}
The final answer is 3 operations.

\textbf{Example 4b: Longest Common Subsequence}
Problem: Find the LCS of ``ABCB'' and ``BDCB''.
State: \texttt{dp[i][j]} is the length of the LCS of the prefixes.
Transition: If match, $1 + \texttt{dp[i-1][j-1]}$. If differ, $\max(\texttt{dp[i-1][j]}, \texttt{dp[i][j-1]})$.

Reconstructing the answer (marked by path traversal):
\begin{enumerate}
    \item Start at bottom-right $(4,4)$ with value 3. Characters \texttt{'B'} and \texttt{'B'} match. Move diagonally up-left to $(3,3)$ with value 2. (Record \texttt{'B'}).
    \item At $(3,3)$ with value 2 (\texttt{'C'} vs \texttt{'C'}), match. Move diagonally up-left to $(2,2)$ with value 1. (Record \texttt{'C'}).
    \item At $(2,2)$ with value 1 (\texttt{'B'} vs \texttt{'D'}), mismatch. Maximum neighbor is left at $(2,1)$ with value 1. Move left to $(2,1)$.
    \item At $(2,1)$ with value 1 (\texttt{'B'} vs \texttt{'B'}), match. Move diagonally up-left to $(1,0)$ with value 0. (Record \texttt{'B'}).
\end{enumerate}
The characters matched backwards are \texttt{'B'}, \texttt{'C'}, \texttt{'B'}. Reversing gives the longest common subsequence: ``BCB''.

\textbf{Example 5: 0/1 Knapsack and Space Optimization}
Problem: Capacity $W=5$. Items (Weight, Value): Item 1 (2, 3), Item 2 (3, 4), Item 3 (4, 8).
Full 2D table trace \texttt{dp[item][w]}:

The 1D version updates a single array. When adding Item 2 (weight 3, value 4), we iterate backwards:

Counter-demonstration of FORWARD loop taking the same item twice:
If we loop forwards for Item 1 (weight 2, value 3):
\texttt{w=2}: \texttt{dp[2]} becomes 3.
\texttt{w=4}: \texttt{dp[4]} checks \texttt{dp[4-2]} which is \texttt{dp[2]}. It sees 3, adds 3, and becomes 6! We just used Item 1 twice. Backwards looping prevents this.

\textbf{Example 6: Equal Partition (Subset-Sum)}
Problem: Array $[1, 5, 11, 5]$. Target sum is $(1+5+11+5)/2 = 11$.
Boolean 1D table (True/False, iterating backwards):

Since \texttt{dp[11]} is True, it is possible.

\textbf{Example 6b: Palindromic Substrings (Interval DP)}
Problem: Count palindromic substrings in ``aba''.
State: \texttt{dp[L][R]} is true if substring from $L$ to $R$ is a palindrome.
\begin{enumerate}
    \item Base case: Length 1 (\texttt{L=R}) is always true (3 palindromes: \texttt{'a'}, \texttt{'b'}, \texttt{'a'}).
    \item Length 2: ``ab'' and ``ba''. Both differ. \texttt{dp[0][1]} = False, \texttt{dp[1][2]} = False.
    \item Length 3: ``aba''. Outer characters \texttt{str[0] == str[2]} (\texttt{'a'} == \texttt{'a'}). Inner string is \texttt{'b'} (length 1, where \texttt{dp[1][1]} is true). Thus \texttt{dp[0][2]} = True (1 palindrome: ``aba'').
\end{enumerate}
Total palindromes: 4.
Table:

\textbf{Example 7: Matrix Chain Multiplication (Interval DP)}
Problem: Dimensions $[10, 30, 5, 60]$. Matrices $A(10 \times 30), B(30 \times 5), C(5 \times 60)$.
Table \texttt{dp[L][R]} holds min cost. Filled by length.
Length 1 (Base): \texttt{dp[0][0]=0, dp[1][1]=0, dp[2][2]=0}.
Length 2:
\texttt{dp[0][1]} (A*B): $10 \times 30 \times 5 = 1500$.
\texttt{dp[1][2]} (B*C): $30 \times 5 \times 60 = 9000$.
Length 3 (A*B*C):
Split at 0: \texttt{dp[0][0] + dp[1][2]} + $10 \times 30 \times 60 = 0 + 9000 + 18000 = 27000$.
Split at 1: \texttt{dp[0][1] + dp[2][2]} + $10 \times 5 \times 60 = 1500 + 0 + 3000 = 4500$.
Min cost is 4500.
Triangular table trace:

\textbf{Example 8: Bitmask Shortest Tour}

Problem: Visit nodes $0, 1, 2$. Distances: $0 \to 1 = 10, 0 \to 2 = 15, 1 \to 2 = 20$.
State: \texttt{dp[mask][last\_visited]}. The mask \texttt{101} in binary (5 in decimal) means nodes 0 and 2 are visited.

\begin{itemize}
    \item Base: \texttt{dp[001][0] = 0}.
    \item Mask 011 (nodes 0,1), last=1: \texttt{dp[011][1] = dp[001][0] + dist[0][1] = 10}.
    \item Mask 101 (nodes 0,2), last=2: \texttt{dp[101][2] = dp[001][0] + dist[0][2] = 15}.
    \item Mask 111 (all nodes):
    \begin{itemize}
        \item End at 1: \texttt{dp[111][1] = dp[101][2] + dist[2][1] = 15 + 20 = 35}.
        \item End at 2: \texttt{dp[111][2] = dp[011][1] + dist[1][2] = 10 + 20 = 30}.
    \end{itemize}
    \item Answer is $\min(35, 30) = \mathbf{30}$.
\end{itemize}

\textbf{Example 9: Solution Reconstruction via Parent Pointers}
Problem: Min coins for amount 6 with $\{1, 3, 4\}$.
We maintain a \texttt{parent} array storing the coin used.

To reconstruct 6: \texttt{parent[6] = 3}. (Amount becomes 6 - 3 = 3).
To reconstruct 3: \texttt{parent[3] = 3}. (Amount becomes 3 - 3 = 0).
Output: 3 + 3.

\section{C++ Implementations}

In modern C++23, we prefer \texttt{std::vector} to raw arrays and use \texttt{constexpr} and \texttt{[[nodiscard]]} where it makes sense.
We define \texttt{INF} carefully to avoid integer overflow if we accidentally add to it. We use \texttt{std::numeric\_limits<int>::max() / 2}. Functions allocating memory (like those returning or creating \texttt{std::vector}) omit \texttt{noexcept}.

\subsection{Tabulated Table}
A tabular approach iteratively fills an array. This explicitly builds the table so you can see all states.

\begin{lstlisting}[caption={Fibonacci with full tabular vector}]
#include <vector>
#include <cstddef>

[[nodiscard]] int fib_tabulated(int n) {
    if (n <= 1) return n;
    // 1. Initialize full DP table
    std::vector<int> dp(n + 1, 0);
    // 2. Set base cases
    dp[0] = 0;
    dp[1] = 1;
    
    // 3. Evaluate strictly from smallest to largest
    for (std::size_t i = 2; i <= static_cast<std::size_t>(n); ++i) {
        dp[i] = dp[i - 1] + dp[i - 2];
    }
    // 4. Return the answer state
    return dp[n];
}
\end{lstlisting}

\subsection{House Robber (O(1) Space)}
Since we only ever need \texttt{dp[i-1]} and \texttt{dp[i-2]}, we can use rolling variables.

\begin{lstlisting}[caption={House Robber optimized to O(1) space}]
#include <vector>
#include <algorithm>
#include <cstddef>

[[nodiscard]] int houseRobber(const std::vector<int>& nums) noexcept {
    if (nums.empty()) return 0;
    if (nums.size() == 1) return nums[0];
    
    // 1. Initialize rolling variables
    int prev2 = nums[0];
    int prev1 = std::max(nums[0], nums[1]);
    
    // 2. Iterate from index 2 onwards
    for (std::size_t i = 2; i < nums.size(); ++i) {
        int current = std::max(prev1, prev2 + nums[i]);
        prev2 = prev1;
        prev1 = current;
    }
    // 3. prev1 holds the final answer
    return prev1;
}
\end{lstlisting}

\subsection{Grid Minimum Path Sum (2D and Rolling 1D)}

\begin{lstlisting}[caption={Grid Minimum Path Sum in 2D and Space Optimized}]
#include <vector>
#include <algorithm>
#include <cstddef>

// 2D Version
[[nodiscard]] int minPathSum2D(const std::vector<std::vector<int>>& grid) {
    std::size_t m = grid.size();
    std::size_t n = grid[0].size();
    std::vector<std::vector<int>> dp = grid;
    
    // 1. Initialize edges (base cases)
    for (std::size_t i = 1; i < m; ++i) dp[i][0] += dp[i - 1][0];
    for (std::size_t j = 1; j < n; ++j) dp[0][j] += dp[0][j - 1];
    
    // 2. Fill the table
    for (std::size_t i = 1; i < m; ++i) {
        for (std::size_t j = 1; j < n; ++j) {
            dp[i][j] += std::min(dp[i - 1][j], dp[i][j - 1]);
        }
    }
    return dp[m - 1][n - 1];
}

// 1D Rolling Array Version
[[nodiscard]] int minPathSum1D(const std::vector<std::vector<int>>& grid) {
    std::size_t m = grid.size();
    std::size_t n = grid[0].size();
    std::vector<int> dp(n, 0);
    
    // 1. Initialize first row
    dp[0] = grid[0][0];
    for (std::size_t j = 1; j < n; ++j) dp[j] = dp[j - 1] + grid[0][j];
    
    // 2. Roll through rows
    for (std::size_t i = 1; i < m; ++i) {
        dp[0] += grid[i][0];
        for (std::size_t j = 1; j < n; ++j) {
            dp[j] = grid[i][j] + std::min(dp[j], dp[j - 1]);
        }
    }
    return dp[n - 1];
}
\end{lstlisting}

\subsection{Longest Common Subsequence with String Reconstruction}

\begin{lstlisting}[caption={LCS with string reconstruction}]
#include <vector>
#include <string>
#include <algorithm>
#include <cstddef>

[[nodiscard]] std::string getLCS(const std::string& s1, const std::string& s2) {
    std::size_t m = s1.length(), n = s2.length();
    std::vector<std::vector<int>> dp(m + 1, std::vector<int>(n + 1, 0));
    
    // 1. Fill the LCS length table
    for (std::size_t i = 1; i <= m; ++i) {
        for (std::size_t j = 1; j <= n; ++j) {
            if (s1[i - 1] == s2[j - 1]) {
                dp[i][j] = dp[i - 1][j - 1] + 1;
            } else {
                dp[i][j] = std::max(dp[i - 1][j], dp[i][j - 1]);
            }
        }
    }
    
    // 2. Reconstruct the string backwards
    std::string lcs = "";
    std::size_t i = m, j = n;
    while (i > 0 && j > 0) {
        if (s1[i - 1] == s2[j - 1]) {
            lcs += s1[i - 1];
            --i; --j;
        } else if (dp[i - 1][j] > dp[i][j - 1]) {
            --i;
        } else {
            --j;
        }
    }
    
    // 3. Reverse the reconstructed path
    std::reverse(lcs.begin(), lcs.end());
    return lcs;
}
\end{lstlisting}

\subsection{Edit Distance}
To find the minimum insertions, deletions, and replacements between two strings, we branch on matching characters or taking the minimum of three single-character operations.

\begin{lstlisting}[caption={Edit Distance using a 2D DP table}]
#include <vector>
#include <string>
#include <algorithm>
#include <cstddef>

[[nodiscard]] int editDistance(const std::string& word1, const std::string& word2) {
    std::size_t m = word1.length();
    std::size_t n = word2.length();
    std::vector<std::vector<int>> dp(m + 1, std::vector<int>(n + 1, 0));

    // 1. Base cases: transforming to/from an empty string
    for (std::size_t i = 0; i <= m; ++i) dp[i][0] = static_cast<int>(i);
    for (std::size_t j = 0; j <= n; ++j) dp[0][j] = static_cast<int>(j);

    // 2. Fill the table
    for (std::size_t i = 1; i <= m; ++i) {
        for (std::size_t j = 1; j <= n; ++j) {
            if (word1[i - 1] == word2[j - 1]) {
                dp[i][j] = dp[i - 1][j - 1]; // Characters match
            } else {
                dp[i][j] = 1 + std::min({
                    dp[i - 1][j],    // Delete from word1
                    dp[i][j - 1],    // Insert into word1
                    dp[i - 1][j - 1] // Replace
                });
            }
        }
    }
    return dp[m][n];
}
\end{lstlisting}

\subsection{Longest Increasing Subsequence ($O(N^2)$)}
At each index $i$, we inspect all preceding elements $j < i$ that are strictly smaller than $\text{nums}[i]$ and take the maximum extending length.

\begin{lstlisting}[caption={Longest Increasing Subsequence in $O(N^2)$ time}]
#include <vector>
#include <algorithm>
#include <cstddef>

[[nodiscard]] int lengthOfLIS(const std::vector<int>& nums) {
    if (nums.empty()) return 0;
    std::size_t n = nums.size();
    std::vector<int> dp(n, 1);
    int max_len = 1;

    // 1. Build LIS ending at index i
    for (std::size_t i = 1; i < n; ++i) {
        for (std::size_t j = 0; j < i; ++j) {
            if (nums[j] < nums[i]) {
                dp[i] = std::max(dp[i], dp[j] + 1);
            }
        }
        max_len = std::max(max_len, dp[i]);
    }
    return max_len;
}
\end{lstlisting}

\subsection{0/1 Knapsack (Value and Weight with 1D Backward Loop)}
Iterating backwards through capacity ensures that each item is added at most once per configuration, preventing unintended duplicate reuse.

\begin{lstlisting}[caption={0/1 Knapsack space-optimized to a 1D backward loop}]
#include <vector>
#include <algorithm>
#include <cstddef>

[[nodiscard]] int knapsack01(int W, const std::vector<int>& weights, const std::vector<int>& values) {
    std::size_t n = weights.size();
    std::vector<int> dp(W + 1, 0);

    // 1. Process each item sequentially
    for (std::size_t i = 0; i < n; ++i) {
        int weight = weights[i];
        int value = values[i];
        // 2. Iterate capacity backwards to use the item at most once
        for (int w = W; w >= weight; --w) {
            dp[w] = std::max(dp[w], dp[w - weight] + value);
        }
    }
    return dp[W];
}
\end{lstlisting}

\subsection{Boolean Knapsack (Subset Sum)}

\begin{lstlisting}[caption={Subset-Sum optimized to 1D boolean array}]
#include <vector>
#include <cstddef>

[[nodiscard]] bool canPartition(const std::vector<int>& nums, int target) {
    // 1. Initialize boolean DP array
    std::vector<bool> dp(target + 1, false);
    dp[0] = true;
    
    // 2. Iterate items, updating capacity BACKWARDS
    for (int num : nums) {
        for (int w = target; w >= num; --w) {
            if (dp[w - num]) {
                dp[w] = true;
            }
        }
    }
    return dp[target];
}
\end{lstlisting}

\subsection{Matrix Chain Multiplication (Interval DP)}

\begin{lstlisting}[caption={Matrix Chain Multiplication by length}]
#include <vector>
#include <algorithm>
#include <limits>
#include <cstddef>

constexpr int INF = std::numeric_limits<int>::max() / 2;

[[nodiscard]] int matrixChain(const std::vector<int>& dims) {
    std::size_t n = dims.size() - 1; // Number of matrices
    std::vector<std::vector<int>> dp(n, std::vector<int>(n, 0));
    
    // 1. Iterate over interval lengths from 2 to n
    for (std::size_t length = 2; length <= n; ++length) {
        for (std::size_t left = 0; left <= n - length; ++left) {
            std::size_t right = left + length - 1;
            dp[left][right] = INF;
            // 2. Try every split point k
            for (std::size_t k = left; k < right; ++k) {
                int cost = dp[left][k] + dp[k + 1][right] + 
                           dims[left] * dims[k + 1] * dims[right + 1];
                dp[left][right] = std::min(dp[left][right], cost);
            }
        }
    }
    return dp[0][n - 1];
}
\end{lstlisting}

\subsection{Bitmask DP (Shortest Tour / Assignment)}
Using bitwise shifts and C++20 \texttt{std::popcount} to process sizes.

\begin{lstlisting}[caption={Shortest tour on a small graph using bitmask DP}]
#include <vector>
#include <algorithm>
#include <limits>
#include <cstddef>
#include <bit> // For std::popcount

[[nodiscard]] int shortestTour(const std::vector<std::vector<int>>& dist) {
    std::size_t n = dist.size();
    std::size_t limit = 1 << n; // 2^N states
    std::vector<std::vector<int>> dp(limit, std::vector<int>(n, INF));
    
    // 1. Base case: starting at node 0
    dp[1][0] = 0;
    
    // 2. Iterate masks
    for (std::size_t mask = 1; mask < limit; ++mask) {
        // Optional: std::popcount(mask) tells us how many nodes are visited
        int visited_count = std::popcount(mask);
        
        for (std::size_t last = 0; last < n; ++last) {
            if (!(mask & (1 << last))) continue; // node 'last' not in mask
            
            // 3. Try to transition to unvisited node 'next'
            for (std::size_t nxt = 0; nxt < n; ++nxt) {
                if (!(mask & (1 << nxt)) && dp[mask][last] != INF) {
                    std::size_t newMask = mask | (1 << nxt);
                    dp[newMask][nxt] = std::min(dp[newMask][nxt], 
                                                dp[mask][last] + dist[last][nxt]);
                }
            }
        }
    }
    
    // 4. Find the minimum tour ending at any node
    int ans = INF;
    for (std::size_t last = 1; last < n; ++last) {
        ans = std::min(ans, dp[limit - 1][last]);
    }
    return ans;
}
\end{lstlisting}

\subsection{Top-down Memoization with Lambda}
We can capture a \texttt{std::vector} in a recursive lambda.

\begin{lstlisting}[caption={Top-down memoization using a recursive lambda}]
#include <vector>
#include <algorithm>
#include <cstddef>

[[nodiscard]] int minPathMemo(const std::vector<std::vector<int>>& grid) {
    std::size_t m = grid.size(), n = grid[0].size();
    std::vector<std::vector<int>> memo(m, std::vector<int>(n, -1));
    
    // 1. Recursive lambda for DP
    auto dfs = [&](this auto& self, std::size_t i, std::size_t j) -> int {
        if (i == 0 && j == 0) return grid[0][0];
        if (i >= m || j >= n) return INF; // Out of bounds safety
        
        // 2. Cache check
        if (memo[i][j] != -1) return memo[i][j];
        
        // 3. Transitions
        int cost_up = (i > 0) ? self(i - 1, j) : INF;
        int cost_left = (j > 0) ? self(i, j - 1) : INF;
        
        // 4. Store and return
        memo[i][j] = grid[i][j] + std::min(cost_up, cost_left);
        return memo[i][j];
    };
    
    return dfs(m - 1, n - 1);
}
\end{lstlisting}

\subsection{Coin Change (Counting Combinations and Path Reconstruction)}
The loop order determines whether we count combinations (outer loop over coins) or permutations (outer loop over amounts). We also demonstrate path reconstruction using parent pointers.

\begin{lstlisting}[caption={Coin change: counting combinations and reconstructing optimal coins}]
#include <vector>
#include <algorithm>
#include <cstddef>

[[nodiscard]] int countCoinWays(const std::vector<int>& coins, int amount) {
    // 1. Initialize combinations array
    std::vector<int> dp(amount + 1, 0);
    dp[0] = 1; // 1 way to make amount 0

    // 2. Outer loop over coins ensures combinations (not permutations)
    for (int coin : coins) {
        for (int a = coin; a <= amount; ++a) {
            dp[a] += dp[a - coin];
        }
    }
    return dp[amount];
}

[[nodiscard]] std::vector<int> getMinCoins(const std::vector<int>& coins, int amount) {
    std::vector<int> dp(amount + 1, INF);
    std::vector<int> parent(amount + 1, -1);
    dp[0] = 0;
    
    // 1. Fill DP and parent pointers
    for (int i = 1; i <= amount; ++i) {
        for (int coin : coins) {
            if (i >= coin && dp[i - coin] != INF) {
                if (dp[i - coin] + 1 < dp[i]) {
                    dp[i] = dp[i - coin] + 1;
                    parent[i] = coin; // Store the coin used
                }
            }
        }
    }
    
    // 2. Reconstruct path
    std::vector<int> result;
    if (dp[amount] == INF) return result;
    
    int curr = amount;
    while (curr > 0) {
        result.push_back(parent[curr]);
        curr -= parent[curr];
    }
    return result;
}
\end{lstlisting}

\section{Problem Pattern Recognition}

Recognizing a DP problem is a skill that requires three steps.

\subsection{Step 1: What does the problem ask?}
\begin{itemize}
    \item The minimum or maximum cost.
    \item The total number of ways (combinations/paths).
    \item Is it possible? (a boolean yes/no).
    \item If the problem asks you to list ALL configurations, it is NOT DP; use Backtracking.
\end{itemize}

\subsection{Step 2: Key Constraint Analysis}
\begin{itemize}
    \item Can you describe the current situation using just a few integer variables (indices)? 
    \item If the state requires you to remember the entire path visited so far, it might be a factorial-time problem unless $N \le 20$ (which implies Bitmask DP).
\end{itemize}

\subsection{Step 3: Threshold Checking}
Check the input constraints. Memory arithmetic rules the state size:
\begin{itemize}
    \item $N \le 5000$: Likely an $O(N^2)$ 1D DP or 2D string DP. $5000 \times 5000$ integers is exactly 100 MB, which fits comfortably, but using rolling rows guarantees safety.
    \item Knapsack capacity $W \times N \le 10^7$: Pseudo-polynomial time is acceptable.
    \item $N \le 20$: Exponential time is small enough. Look for bitmask DP.
    \item If both $N$ and the values are huge, it is probably not DP (or it is a Greedy problem).
\end{itemize}

\noindent
\vspace{0.3cm}
\begin{tabularx}{\textwidth}{|p{4.2cm}|X|p{4cm}|}
\hline
\textbf{Keywords in the Problem} & \textbf{Plain English Meaning} & \textbf{Tool to Use} \\ \hline
number of ways to reach & Counting paths or combinations & DP (add transitions) \\ \hline
number of distinct subsequences / ways to form a string & Counting matches between characters & DP on Strings \\ \hline
minimum cost/minimum number of ... & Shortest path in state graph & DP (min transition) \\ \hline
minimum number of coins / steps to reach a target & Shortest path bounded by transitions & 1D Coin Change \\ \hline
maximum value without exceeding capacity & Choose items with a budget & 0/1 Knapsack \\ \hline
maximum sum of a non-adjacent selection & Pick separated elements & House Robber (1D) \\ \hline
longest common subsequence & Match characters in order & Two-Sequence DP \\ \hline
longest palindromic subsequence/substring & Palindrome check in a window & Interval DP (expand) \\ \hline
minimum edits to transform & Insert/Delete/Replace ops & Edit Distance DP \\ \hline
longest increasing subsequence & Pick elements strictly increasing & 1D DP (LIS) \\ \hline
can we split into equal sum halves & Target sum is exactly half total & Subset-Sum (Knapsack) \\ \hline
partition into k groups with minimum cost & Cost grouped by elements & DP (Index and K) \\ \hline
minimum cost to merge/multiply a sequence & Cost depends on final merge point & Interval DP \\ \hline
stock-style buy/sell with states (holding / not holding) & State transitions with multiple roles & State-Machine DP \\ \hline
unbounded reuse of items & Take the same item infinitely & Unbounded Knapsack (forward loop) \\ \hline
visit every node exactly once, $n \le 20$ & Traveling Salesman style & Bitmask DP \\ \hline
count ways modulo & Answers get too large & DP (apply \% at adds) \\ \hline
grid paths with obstacles & Navigating 2D array & 2D Grid DP \\ \hline
no two adjacent chosen & House robber constraint & 1D DP \\ \hline
\end{tabularx}
\vspace{0.3cm}

\section{Summary and Complexity Reference}

\noindent
\begin{tabularx}{\textwidth}{|l|l|l|X|}
\hline
\textbf{Task / Operation} & \textbf{Time} & \textbf{Space} & \textbf{Use When} \\ \hline
1D DP (Stairs, Robber) & $O(N)$ & $O(N)$ & Only depends on index \\ \hline
1D DP Space Optimized & $O(N)$ & $O(1)$ & Depends only on $k$ last states \\ \hline
2D Grid DP & $O(N \times M)$ & $O(N \times M)$ & Moving in a matrix \\ \hline
2D Space Optimized & $O(N \times M)$ & $O(\min(N, M))$ & State only reads previous row \\ \hline
0/1 Knapsack & $O(N \times W)$ & $O(W)$ & Maximize value with weight limit \\ \hline
Longest Common Subsequence & $O(N \times M)$ & $O(N \times M)$ & Subsequence match two strings \\ \hline
Longest Increasing Subsequence & $O(N^2)$ & $O(N)$ & Standard LIS DP formulation \\ \hline
Interval DP & $O(N^3)$ & $O(N^2)$ & Merging adjacent items \\ \hline
Bitmask DP & $O(2^N \times N)$ & $O(2^N \times N)$ & Set of visited items, $N \le 20$ \\ \hline
\end{tabularx}

\section{Glossary}

\subsection{Core Concepts}
\begin{itemize}
    \item \textbf{0/1 Knapsack}: A constraint where an item can be chosen at most once. Represented by a backward capacity loop.
    \item \textbf{Answer State}: The specific state holding the final answer to the overall problem.
    \item \textbf{Backtracking}: A brute-force search that explores all valid permutations. Used when you must list all configurations.
    \item \textbf{Base Case}: The smallest instances requiring no recursive calculation.
    \item \textbf{Bitmask}: Using the binary representation of an integer to track subsets (for example, bit 2 is 1 if item 2 is used).
    \item \textbf{Directed Acyclic Graph (DAG)}: A graph with one-way edges and no cycles. DP states form a DAG.
    \item \textbf{Edit Distance (Levenshtein)}: Minimum operations (insert, delete, replace) to transform one sequence into another.
    \item \textbf{Evaluation Order}: The sequence in which subproblems must be solved.
    \item \textbf{Greedy}: An algorithm that picks the locally best choice. Often fails when future choices depend on the current one.
    \item \textbf{Longest Common Subsequence (LCS)}: The longest shared sequence of characters between two strings in order.
    \item \textbf{Longest Increasing Subsequence (LIS)}: The longest strictly increasing subsequence of an array.
    \item \textbf{Memoization}: A top-down optimization caching recursive function results.
    \item \textbf{Optimal Substructure}: A property where the optimal solution contains optimal solutions to its subproblems.
    \item \textbf{Overlapping Subproblems}: When an algorithm solves the exact same subproblem repeatedly.
    \item \textbf{Parent Pointer}: An array used during reconstruction to explicitly store the decision made at a state.
    \item \textbf{Popcount}: The count of set bits (1s) in a binary representation.
    \item \textbf{Pseudo-polynomial}: Complexity dependent on the numeric value of the input rather than just the number of items.
    \item \textbf{Reconstruction}: Tracing backwards from the final state to deduce the exact sequence of optimal choices.
    \item \textbf{Rolling Array}: A space optimization keeping only the required slice of previous computations.
    \item \textbf{State}: A distinct subproblem configuration.
    \item \textbf{Subsequence vs Substring}: A subsequence is derived by deleting elements; a substring must be contiguous.
    \item \textbf{Tabulation}: A bottom-up technique using iterative loops to fill a table.
    \item \textbf{Topological Order}: An ordering of DAG nodes such that every node appears after all its dependencies.
    \item \textbf{Transition (Recurrence)}: The rule to calculate a state from smaller states.
    \item \textbf{Unbounded Knapsack}: Items can be chosen infinitely many times. Represented by a forward capacity loop.
\end{itemize}

\subsection{Mathematical Symbols}
\begin{itemize}
    \item \textbf{$2^N$}: Exponential growth. Characteristic of subsets and Bitmask DP limits.
    \item \textbf{$\le$}: Less than or equal to.
    \item \textbf{$N!$}: Factorial growth. Characteristic of permutation generation.
    \item \textbf{$O(N)$ (Big-O)}: Describes the upper bound of an algorithm's complexity.
    \item \textbf{$\max / \min$}: Functions returning the largest or smallest value.
    \item \textbf{$\min_k$}: The minimum over all possible values of variable $k$. Used in Interval DP.
    \item \textbf{\texttt{dp[i][j]}}: Array access notation for the stored answer of a 2D state.
    \item \textbf{$\implies$}: Logical implication (if A is true, then B is true).
    \item \textbf{$\sum$}: Summation of a sequence of values.
    \item \textbf{INF (Sentinel)}: A large constant representing infinity.
    \item \textbf{$\%$ MOD}: Modulo arithmetic keeping values within bounds.
    \item \textbf{\texttt{<<}, \texttt{\&}, \texttt{|}}: Bitwise left-shift, AND, and OR operations used for manipulating Bitmasks.
\end{itemize}

\section{Further Reading}

\subsection*{Free Online Resources (Start Here)}
\begin{itemize}
    \item \textbf{cp-algorithms.com}: Dynamic Programming pages (Highly recommended free resource for exact code templates).
    \item \textbf{Project Euler}: Practice problems. Problem 15 (Lattice Paths), Problem 31 (Coin Sums), Problems 18 and 67 (Maximum path sum in a triangle), Problems 76 and 77 (Partitions).
    \item \textbf{Art of Problem Solving (AoPS) Wiki}: Excellent theoretical breakdown of combinations and path counting.
    \item \textbf{Brilliant.org}: Visual guides on dynamic programming transitions.
    \item \textbf{Wikipedia}: Articles on ``Dynamic Programming'', ``Knapsack Problem'', ``Longest Common Subsequence'', and ``Levenshtein Distance''.
\end{itemize}

\subsection*{Books (When You Are Ready to Go Deeper)}
\begin{itemize}
    \item \textit{The Algorithm Design Manual} by Steven Skiena: Intermediate. Great real-world intuition for DP.
    \item \textit{Competitive Programming} by Steven Halim: Intermediate. Actionable guidance for translating math constraints into code.
    \item \textit{Algorithms} by Dasgupta, Papadimitriou, and Vazirani: Friendly intermediate math textbook with excellent DP visual examples.
    \item \textit{Introduction to Algorithms} (CLRS): Advanced. Detailed proofs for optimal substructure properties.
\end{itemize}

\section*{Version History}
\addcontentsline{toc}{section}{Version History}

\noindent
\begin{tabularx}{\textwidth}{|l|X|}
\hline
\textbf{Date} & \textbf{Change} \\ \hline
2026-10-03 & Document created. \\ \hline
\end{tabularx}

\end{document}

